% 第19章习题解答
% 顶层题目 42 道；按明确小问计 53 个作答单元。

\chapter*{第19章：偏导数与全微分习题解答}
\addcontentsline{toc}{chapter}{第19章：偏导数与全微分习题解答}
\section*{19.2.3　思考题}

\subsection*{思考题 1}
\begin{xsolution}
$f_x(x_0,y_0)$ 存在，按偏导数定义，就是一元函数
\[
  \varphi(x)=f(x,y_0)
\]
在 $x_0$ 可导。一元函数可导必连续，因此
\[
  \lim_{x\to x_0}f(x,y_0)=f(x_0,y_0).
\]

但不能由此断言当 $y_1$ 充分接近 $y_0$ 时，$x\mapsto f(x,y_1)$ 在 $x_0$ 连续。例如取 $(x_0,y_0)=(0,0)$，定义
\[
  f(x,y)=
  \begin{cases}
    1,&x\ne0,\ y\ne0,\\
    0,&x=0\text{ 或 }y=0.
  \end{cases}
\]
则 $f(x,0)\equiv0$，故 $f_x(0,0)=0$。然而对任意 $y_1\ne0$，都有
\[
  f(0,y_1)=0,\qquad f(x,y_1)=1\quad(x\ne0),
\]
所以 $x\mapsto f(x,y_1)$ 在 $x=0$ 不连续。由于这样的 $y_1$ 可以任意接近 $0$，所述进一步断言不成立。
\end{xsolution}

\subsection*{思考题 2}
\begin{xsolution}
设 $f$ 在 $(x_0,y_0)$ 可微，则
\[
  \Delta f=A\Delta x+B\Delta y+\littleo(r),
  \qquad r=\sqrt{(\Delta x)^2+(\Delta y)^2}.
\]

先证明性质 (1)。令 $\Delta y=0$，则 $r=|\Delta x|$，于是
\[
  \frac{f(x_0+\Delta x,y_0)-f(x_0,y_0)}{\Delta x}
  =A+\frac{\littleo(|\Delta x|)}{\Delta x}.
\]
令 $\Delta x\to0$，右端第二项趋于 $0$，故
\[
  f_x(x_0,y_0)=A.
\]
同理令 $\Delta x=0$，得到
\[
  f_y(x_0,y_0)=B.
\]
因此
\[
  \dd f=f_x(x_0,y_0)\dd x+f_y(x_0,y_0)\dd y.
\]

再证明性质 (2)。由可微公式
\[
  |\Delta f|
  \le |A|\,|\Delta x|+|B|\,|\Delta y|+|\littleo(r)|
  \le (|A|+|B|)r+|\littleo(r)|.
\]
当 $r\to0$ 时右端趋于 $0$，故 $\Delta f\to0$，即 $f$ 在 $(x_0,y_0)$ 连续。
\end{xsolution}

\subsection*{思考题 3}
\begin{xsolution}
分别给出反例。

\textbf{(1)} 定义
\[
  f(x,y)=
  \begin{cases}
    \dfrac{xy}{x^2+y^2},&(x,y)\ne(0,0),\\
    0,&(x,y)=(0,0).
  \end{cases}
\]
在原点以外 $f$ 光滑；而在原点
\[
  f_x(0,0)=\lim_{h\to0}\frac{f(h,0)}h=0,
  \qquad
  f_y(0,0)=\lim_{h\to0}\frac{f(0,h)}h=0.
\]
故两个偏导数在原点的一个邻域内处处存在。但沿 $y=x\ne0$ 有
\[
  f(x,x)=\frac12,
\]
故 $f$ 在原点不连续，从而也不可微。

\textbf{(2)} 取
\[
  f(x,y)=\sqrt{x^2+y^2}.
\]
该函数在原点连续，但
\[
  \frac{f(h,0)-f(0,0)}h=\frac{|h|}{h}
\]
在 $h\to0$ 时无极限，所以 $f_x(0,0)$ 不存在，因而 $f$ 在原点不可微。

\textbf{(3)} 定义
\[
  f(x,y)=
  \begin{cases}
    xy\sin\dfrac1{x^2+y^2},&(x,y)\ne(0,0),\\
    0,&(x,y)=(0,0).
  \end{cases}
\]
在原点 $f_x(0,0)=f_y(0,0)=0$。又令 $r=\sqrt{x^2+y^2}$，则
\[
  \frac{|f(x,y)|}{r}
  \le\frac{|xy|}{r}
  \le\frac{r}{2}\to0,
\]
故 $f$ 在原点可微，且全微分为零。

但当 $(x,y)\ne(0,0)$ 时
\[
  f_x=y\sin\frac1{x^2+y^2}
  -\frac{2x^2y}{(x^2+y^2)^2}\cos\frac1{x^2+y^2}.
\]
沿 $x=y=t$，取 $t_k=(4\pi k)^{-1/2}$，则
\[
  \frac1{2t_k^2}=2\pi k,
\]
从而
\[
  f_x(t_k,t_k)=-\frac1{2t_k}\to-\infty.
\]
所以 $f_x$ 在原点不连续。类似地 $f_y$ 在原点也不连续。这说明可微并不推出偏导数连续。
\end{xsolution}

\subsection*{思考题 4}
\begin{xsolution}
记 $p=(x_0,y_0)$，并设 $f_x(p)$ 存在、$f_y$ 在 $p$ 连续。将增量分解为
\begin{align*}
  \Delta f
  &=f(x_0+\Delta x,y_0+\Delta y)-f(x_0+\Delta x,y_0)\\
  &\quad+f(x_0+\Delta x,y_0)-f(x_0,y_0).
\end{align*}

由于 $f_y$ 在 $p$ 连续，因而在 $p$ 的某邻域内存在。对第一项把 $x=x_0+\Delta x$ 固定，对变量 $y$ 应用一元函数中值定理，得
\[
  f(x_0+\Delta x,y_0+\Delta y)-f(x_0+\Delta x,y_0)
  =f_y(x_0+\Delta x,y_0+\theta\Delta y)\Delta y,
\]
其中 $0<\theta<1$。由 $f_y$ 在 $p$ 连续，
\[
  f_y(x_0+\Delta x,y_0+\theta\Delta y)
  =f_y(p)+\littleo(1),
\]
故第一项等于
\[
  f_y(p)\Delta y+\littleo(r),
  \qquad r=\sqrt{(\Delta x)^2+(\Delta y)^2}.
\]
另一方面，由 $f_x(p)$ 存在，
\[
  f(x_0+\Delta x,y_0)-f(x_0,y_0)
  =f_x(p)\Delta x+\littleo(|\Delta x|)
  =f_x(p)\Delta x+\littleo(r).
\]
两式相加得到
\[
  \Delta f=f_x(p)\Delta x+f_y(p)\Delta y+\littleo(r),
\]
所以 $f$ 在 $p$ 可微。
\end{xsolution}

\subsection*{思考题 5}
\begin{xsolution}
应取常数值。因为 $\dd z\equiv0$，由可微函数全微分的表示式知
\[
  f_x(x,y)=f_y(x,y)=0,
  \qquad (x,y)\in D.
\]
任取 $(x_1,y_1),(x_2,y_2)\in D=(a,b)\times(c,d)$。矩形区域保证线段
\[
  \{(x,y_1):x_1\le x\le x_2\},
  \qquad
  \{(x_2,y):y_1\le y\le y_2\}
\]
均在 $D$ 内。固定 $y=y_1$，一元函数 $x\mapsto f(x,y_1)$ 的导数恒为零，故
\[
  f(x_1,y_1)=f(x_2,y_1).
\]
再固定 $x=x_2$，同理由 $f_y=0$ 得
\[
  f(x_2,y_1)=f(x_2,y_2).
\]
因此
\[
  f(x_1,y_1)=f(x_2,y_2).
\]
两点任意，所以 $f$ 在 $D$ 上为常值函数。
\end{xsolution}

\section*{19.2.4　练习题}

\subsection*{练习题 1}
\begin{xsolution}
记 $r=\sqrt{x^2+y^2}$。

\textbf{(1)} 当 $(x,y)\ne(0,0)$ 时，利用 $|\sin r^2|\le r^2$，有
\[
  |f(x,y)|
  =\frac{\sqrt{|x-y|}}{r^2}|\sin r^2|
  \le\sqrt{|x-y|}.
\]
又
\[
  |x-y|\le |x|+|y|\le\sqrt2\,r,
\]
故
\[
  |f(x,y)|\le 2^{1/4}r^{1/2}\to0=f(0,0).
\]
所以 $f$ 在 $(0,0)$ 连续。

\textbf{(2)} 取 $y=0$。当 $h\ne0$ 时
\[
  f(h,0)=\frac{\sqrt{|h|}}{h^2}\sin h^2,
\]
因而
\[
  \frac{f(h,0)-f(0,0)}h
  =\frac{\sqrt{|h|}\sin h^2}{h^3}
  =\frac{\sin h^2}{h^2}\cdot\frac{\sqrt{|h|}}h.
\]
第一因子趋于 $1$，第二因子的绝对值为 $|h|^{-1/2}\to+\infty$，故 $f_x(0,0)$ 不存在。可微必有偏导数，因此 $f$ 在 $(0,0)$ 不可微。
\end{xsolution}

\subsection*{练习题 2}
\begin{xsolution}
\textbf{(1)} 沿坐标轴有 $f(h,0)=f(0,h)=0$，故
\[
  f_x(0,0)=0,
  \qquad
  f_y(0,0)=0.
\]

\textbf{(2)} 当 $(x,y)\ne(0,0)$ 时
\[
  f_x(x,y)=y\sin\frac1{x^2+y^2}
  -\frac{2x^2y}{(x^2+y^2)^2}\cos\frac1{x^2+y^2},
\]
\[
  f_y(x,y)=x\sin\frac1{x^2+y^2}
  -\frac{2xy^2}{(x^2+y^2)^2}\cos\frac1{x^2+y^2}.
\]
令 $x=y=t$，则
\[
  f_x(t,t)=f_y(t,t)
  =t\sin\frac1{2t^2}-\frac1{2t}\cos\frac1{2t^2}.
\]
取 $t_k=(4\pi k)^{-1/2}$，有
\[
  f_x(t_k,t_k)=f_y(t_k,t_k)=-\frac1{2t_k},
\]
其绝对值趋于无穷。因此 $f_x,f_y$ 在 $(0,0)$ 都不连续。

\textbf{(3)} 在原点线性主部为零。由
\[
  |f(x,y)|\le|xy|\le\frac{x^2+y^2}{2}=\frac{r^2}{2}
\]
可得
\[
  \frac{|f(x,y)-f(0,0)-f_x(0,0)x-f_y(0,0)y|}{r}
  \le\frac r2\to0.
\]
故 $f$ 在 $(0,0)$ 可微，且 $\dd f(0,0)=0$。
\end{xsolution}

\subsection*{练习题 3}
\begin{xsolution}
任取 $p=(x_0,y_0,z_0)\in D$。因为 $D$ 为开集，可取 $\rho>0$，使以 $p$ 为中心的一个小长方体完全包含于 $D$。设在 $D$ 上
\[
  |f_x|\le M,
  \qquad
  |f_y|\le M
\]
（若两个界不同，取其较大者）。对充分接近 $p$ 的 $(x,y,z)$，分解
\begin{align*}
  f(x,y,z)-f(x_0,y_0,z_0)
  &=f(x,y,z)-f(x_0,y,z)\\
  &\quad+f(x_0,y,z)-f(x_0,y_0,z)\\
  &\quad+f(x_0,y_0,z)-f(x_0,y_0,z_0).
\end{align*}
前两项分别对 $x,y$ 应用一元函数中值定理，得到
\[
  |f(x,y,z)-f(x_0,y,z)|\le M|x-x_0|,
\]
\[
  |f(x_0,y,z)-f(x_0,y_0,z)|\le M|y-y_0|.
\]
对固定的 $(x_0,y_0)$，题设说明 $z\mapsto f(x_0,y_0,z)$ 在 $z_0$ 连续，因此第三项在 $z\to z_0$ 时趋于 $0$。于是当 $(x,y,z)\to p$ 时三项之和趋于 $0$，故 $f$ 在 $p$ 连续。由于 $p$ 任意，$f$ 在 $D$ 上连续。
\end{xsolution}

\subsection*{练习题 4}
\begin{xsolution}
设
\[
  u(x,y)=\ln(1+x^2+y^2).
\]
则
\[
  u_x=\frac{2x}{1+x^2+y^2},
  \qquad
  u_y=\frac{2y}{1+x^2+y^2}.
\]
在 $(1,2)$ 处
\[
  u_x(1,2)=\frac13,
  \qquad
  u_y(1,2)=\frac23.
\]
所以
\[
  \boxed{\dd u\big|_{(1,2)}=\frac13\dd x+\frac23\dd y}.
\]
\end{xsolution}

\subsection*{练习题 5}
\begin{xsolution}
体积
\[
  V(D,H)=\frac\pi4D^2H.
\]
其全微分为
\[
  \dd V=\frac\pi2DH\,\dd D+\frac\pi4D^2\,\dd H.
\]
故绝对误差可用全微分估计为
\begin{align*}
  |\Delta V|
  &\approx |\dd V|\\
  &\le \frac\pi2D_0H_0|\Delta D|
       +\frac\pi4D_0^2|\Delta H|\\
  &\le \frac\pi2(10.44)(18.36)(0.02)
       +\frac\pi4(10.44)^2(0.01)\\
  &=2.189268\pi\approx6.88.
\end{align*}

又
\[
  \frac{\dd V}{V}=2\frac{\dd D}{D}+\frac{\dd H}{H},
\]
所以相对误差满足
\begin{align*}
  \left|\frac{\Delta V}{V}\right|
  &\approx\left|\frac{\dd V}{V}\right|\\
  &\le2\frac{0.02}{10.44}+\frac{0.01}{18.36}\\
  &\approx0.004376,
\end{align*}
即约为
\[
  \boxed{0.438\%}.
\]
相应的绝对误差估计为
\[
  \boxed{|\Delta V|\lesssim6.88}.
\]
\end{xsolution}

\subsection*{练习题 6}
\begin{xsolution}
因为 $f_y$ 在紧矩形
\[
  I=[a,b]\times[c,d]
\]
上连续，所以 $f_y$ 在 $I$ 上有界。令
\[
  M=\max_{(x,y)\in I}|f_y(x,y)|<\infty.
\]
固定任意 $x\in[a,b]$，对一元函数 $y\mapsto f(x,y)$ 在 $y_1,y_2$ 之间应用中值定理，存在介于 $y_1,y_2$ 之间的 $\xi$，使
\[
  f(x,y_1)-f(x,y_2)=f_y(x,\xi)(y_1-y_2).
\]
于是
\[
  |f(x,y_1)-f(x,y_2)|
  \le M|y_1-y_2|.
\]
若 $M>0$，取 $L=M$；若 $M=0$，可任取 $L>0$。这个 $L$ 与 $x$ 无关，故 $f$ 对 $y$ 满足一致 Lipschitz 条件。
\end{xsolution}

\subsection*{练习题 7}
\begin{xsolution}
任取 $p\in D$。因为 $D$ 是区域，特别是开集，所以存在以 $p$ 为中心的小矩形 $R\subset D$。在 $R$ 内任取两点 $(x_1,y_1),(x_2,y_2)$。先固定 $y=y_1$，由 $f_x=0$ 和一元函数中值定理，
\[
  f(x_1,y_1)=f(x_2,y_1).
\]
再固定 $x=x_2$，由 $f_y=0$，
\[
  f(x_2,y_1)=f(x_2,y_2).
\]
所以 $f$ 在每个这样的 $R$ 上为常数，即 $f$ 在 $D$ 上局部常值。

局部常值函数在每个连通分支上为常数，而“区域”按定义是连通开集。因此 $f$ 在整个 $D$ 上为常值函数。
\end{xsolution}

\subsection*{练习题 8}
\begin{xsolution}
由
\[
  u^2+v^2=C
\]
分别对 $x,y$ 求偏导，得
\[
  u u_x+v v_x=0,
  \qquad
  u u_y+v v_y=0.
\]
再利用
\[
  u_x=v_y,
  \qquad
  u_y=-v_x,
\]
可改写为
\[
  u u_x-v u_y=0,
  \qquad
  v u_x+u u_y=0.
\]
即
\[
  \begin{pmatrix}
    u&-v\\
    v&u
  \end{pmatrix}
  \binom{u_x}{u_y}=0.
\]
其系数行列式为
\[
  u^2+v^2=C.
\]
若 $C=0$，则 $u=v=0$，结论显然。若 $C>0$，系数矩阵可逆，从而
\[
  u_x=u_y=0.
\]
再由 Cauchy--Riemann 型关系得到
\[
  v_x=-u_y=0,
  \qquad
  v_y=u_x=0.
\]
由上一题的结论，$u,v$ 在区域 $\Omega$ 上都为常值函数。
\end{xsolution}

\subsection*{练习题 9}
\begin{xsolution}
设
\[
  |f_y(x,y)|\le M,
  \qquad (x,y)\in G.
\]
当 $(x,y)$ 充分接近 $(0,0)$ 时，竖直线段从 $(x,0)$ 到 $(x,y)$ 完全包含在圆盘 $G$ 中。对 $y$ 应用一元函数中值定理，得
\[
  |f(x,y)-f(x,0)|\le M|y|.
\]
于是
\[
  |f(x,y)-f(0,0)|
  \le M|y|+|f(x,0)-f(0,0)|.
\]
当 $(x,y)\to(0,0)$ 时，第一项趋于 $0$；第二项由题设 $f(x,0)$ 在 $x=0$ 连续也趋于 $0$。因此
\[
  \lim_{(x,y)\to(0,0)}f(x,y)=f(0,0),
\]
即 $f$ 在 $(0,0)$ 连续。
\end{xsolution}

\section*{19.3.4　练习题}

\subsection*{练习题 1}
\begin{xsolution}
由链式法则
\[
  \frac{\partial u}{\partial t}
  =\ee^x\frac{\partial x}{\partial t}
   +\cos y\frac{\partial y}{\partial t}+1.
\]
因为 $x=st,y=s+t$，故
\[
  \frac{\partial x}{\partial t}=s,
  \qquad
  \frac{\partial y}{\partial t}=1.
\]
因此
\[
  \boxed{\frac{\partial u}{\partial t}
  =s\ee^{st}+\cos(s+t)+1}.
\]
\end{xsolution}

\subsection*{练习题 2}
\begin{xsolution}
由于 $s=x/y$、$t=y/z$，以下计算在 $y\ne0$、$z\ne0$ 的定义域内进行。记 $f_s=\frac{\partial f}{\partial s}$，$f_t=\frac{\partial f}{\partial t}$，它们均在
\[
  s=\frac xy,
  \qquad
  t=\frac yz
\]
处取值。由链式法则
\[
  \boxed{\frac{\partial u}{\partial x}=\frac1y f_s},
\]
\[
  \boxed{\frac{\partial u}{\partial y}
  =-\frac{x}{y^2}f_s+\frac1z f_t},
\]
\[
  \boxed{\frac{\partial u}{\partial z}
  =-\frac{y}{z^2}f_t}.
\]
\end{xsolution}

\subsection*{练习题 3}
\begin{xsolution}
把 $f$ 的三个自变量依次记为 $(s,t,y)$。由链式法则
\[
  \boxed{
  \frac{\partial u}{\partial x}
  =f_s\varphi_x+f_t\psi_x},
\]
而对 $y$ 求偏导时，第三个自变量 $y$ 本身还贡献一项，因此
\[
  \boxed{
  \frac{\partial u}{\partial y}
  =f_s\varphi_y+f_t\psi_y+f_y}.
\]
其中 $f_s,f_t,f_y$ 均在
\[
  (s,t,y)=(\varphi(x,y),\psi(x,y),y)
\]
处取值。
\end{xsolution}

\subsection*{练习题 4}
\begin{xsolution}
由
\[
  u=\ee^{2st}\sin(t+s^2)
\]
直接求偏导，得
\[
  \boxed{
  \frac{\partial u}{\partial s}
  =\ee^{2st}\left[2t\sin(t+s^2)+2s\cos(t+s^2)\right]},
\]
\[
  \boxed{
  \frac{\partial u}{\partial t}
  =\ee^{2st}\left[2s\sin(t+s^2)+\cos(t+s^2)\right]}.
\]
\end{xsolution}

\subsection*{练习题 5}
\begin{xsolution}
令
\[
  q=ax^2+by^2+cz^2.
\]
则 $u=f(q)$，故
\[
  \dd u=f'(q)\dd q
  =f'(q)(2ax\dd x+2by\dd y+2cz\dd z).
\]
即
\[
  \boxed{\dd u
  =2f'(ax^2+by^2+cz^2)
  (ax\dd x+by\dd y+cz\dd z)}.
\]
\end{xsolution}

\subsection*{练习题 6}
\begin{xsolution}
在 $(1,1,1)$ 的邻域内可写
\[
  f(x,y,z)=\exp\left[\frac1z(\ln x-\ln y)\right].
\]
因此
\[
  f_x=\frac{f}{zx},
  \qquad
  f_y=-\frac{f}{zy},
  \qquad
  f_z=-\frac{f}{z^2}(\ln x-\ln y).
\]
在 $(1,1,1)$ 处，$f=1$，于是
\[
  f_x=1,
  \qquad f_y=-1,
  \qquad f_z=0.
\]
所以
\[
  \boxed{\dd f(1,1,1)=\dd x-\dd y}.
\]
\end{xsolution}

\subsection*{练习题 7}
\begin{xsolution}
令
\[
  p=xy,
  \qquad q=yz,
  \qquad w=F(p,q).
\]
则
\[
  w_x=F_p y,
  \qquad
  w_z=F_q y,
  \qquad
  w_y=F_p x+F_q z.
\]
从而
\[
  xw_x+zw_z
  =xyF_p+zyF_q
  =y(xF_p+zF_q)
  =yw_y.
\]
故
\[
  \boxed{x\frac{\partial w}{\partial x}
  +z\frac{\partial w}{\partial z}
  =y\frac{\partial w}{\partial y}}.
\]
\end{xsolution}

\subsection*{练习题 8}
\begin{xsolution}
设 $s=xy$。由链式法则
\[
  z_x=f'(s)y,
  \qquad
  z_y=f'(s)x.
\]
故
\[
  xz_x-yz_y
  =xyf'(s)-xyf'(s)=0.
\]
即
\[
  \boxed{x\frac{\partial z}{\partial x}
  -y\frac{\partial z}{\partial y}=0}.
\]
\end{xsolution}

\subsection*{练习题 9}
\begin{xsolution}
令
\[
  U(r,\theta)=F(r\cos\theta,r\sin\theta).
\]
当 $r>0$ 时，由链式法则
\begin{align*}
  \frac{\partial U}{\partial r}
  &=F_x\cos\theta+F_y\sin\theta\\
  &=\frac1r(xF_x+yF_y)=0.
\end{align*}
因此对每个固定的 $\theta$，$U(r,\theta)$ 与 $r$ 无关，故存在一元函数 $h$，使
\[
  \boxed{F(r\cos\theta,r\sin\theta)=h(\theta)}
\]
（在 $r>0$ 的定义域内）。也就是说，$F$ 在极坐标下只是 $\theta$ 的函数。
\end{xsolution}

\subsection*{练习题 10}
\begin{xsolution}
在 $F=f(x)g(y)$ 非零的连通部分上，由题设 $F(r\cos\theta,r\sin\theta)=h(\theta)$ 可知 $F$ 与 $r$ 无关。因此
\[
  0=\frac{\partial F}{\partial r}
  =F_x\cos\theta+F_y\sin\theta.
\]
乘以 $r$，得到
\[
  xf'(x)g(y)+yf(x)g'(y)=0.
\]
在 $x,y,f(x),g(y)$ 均非零处除以 $f(x)g(y)$，有
\[
  x\frac{f'(x)}{f(x)}
  =-y\frac{g'(y)}{g(y)}.
\]
左端只依赖 $x$，右端只依赖 $y$，故二者必为同一常数 $\lambda$：
\[
  x\frac{f'(x)}{f(x)}=\lambda,
  \qquad
  y\frac{g'(y)}{g(y)}=-\lambda.
\]
于是
\[
  f(x)=C_1|x|^\lambda,
  \qquad
  g(y)=C_2|y|^{-\lambda}
\]
（在固定符号的区间内）。故在每个不穿过坐标轴的连通分支上
\[
  \boxed{F(x,y)=C\left|\frac{x}{y}\right|^\lambda}.
\]
若限定在 $x>0,y>0$，则可写成
\[
  \boxed{F(x,y)=C\left(\frac{x}{y}\right)^\lambda}.
\]
恒零函数也是题设的解。若还要求跨坐标轴或在原点连续可微，则应再按相应定义域的正则性要求对上述常数和指数作限制。
\end{xsolution}

\subsection*{练习题 11}
\begin{xsolution}
令
\[
  \xi=x+at,
  \qquad
  \eta=x-at,
  \qquad
  u=\varphi(\xi)+\psi(\eta).
\]
则
\[
  u_x=\varphi'(\xi)+\psi'(\eta),
  \qquad
  u_t=a\varphi'(\xi)-a\psi'(\eta).
\]
再求一次偏导：
\[
  u_{xx}=\varphi''(\xi)+\psi''(\eta),
\]
\[
  u_{tt}=a^2\varphi''(\xi)+a^2\psi''(\eta).
\]
所以
\[
  \boxed{u_{tt}-a^2u_{xx}=0}.
\]
\end{xsolution}

\subsection*{练习题 12}
\begin{xsolution}
为书写简洁，记复合函数仍为
\[
  w(u,v)=w(x(u,v),y(u,v)).
\]
题设关系为
\[
  x_u=y_v,
  \qquad
  x_v=-y_u.
\]
题目涉及关于 $u,v$ 的二阶偏导，因此以下按这些复合函数所需的二阶偏导存在且混合偏导可交换来计算。于是由题设关系进一步有
\[
  x_{uu}+x_{vv}=y_{vu}-y_{uv}=0,
\]
\[
  y_{uu}+y_{vv}=-x_{vu}+x_{uv}=0.
\]
此外
\[
  x_u y_u+x_v y_v
  =y_vy_u-y_uy_v=0,
\]
以及
\[
  x_u^2+x_v^2=y_u^2+y_v^2.
\]

\textbf{(1)} 由二阶链式法则
\begin{align*}
  w_{uu}+w_{vv}
  &=(x_u^2+x_v^2)w_{xx}
    +2(x_uy_u+x_vy_v)w_{xy}\\
  &\quad +(y_u^2+y_v^2)w_{yy}
    +(x_{uu}+x_{vv})w_x
    +(y_{uu}+y_{vv})w_y\\
  &=(x_u^2+x_v^2)(w_{xx}+w_{yy})=0.
\end{align*}
故
\[
  \boxed{w_{uu}+w_{vv}=0}.
\]

\textbf{(2)} 把上一结论应用于 $w(x,y)=xy$。因为
\[
  w_{xx}=w_{yy}=0,
\]
所以 $w=xy$ 是关于 $(x,y)$ 的调和函数。故复合后
\[
  \boxed{
  \frac{\partial^2(xy)}{\partial u^2}
  +\frac{\partial^2(xy)}{\partial v^2}=0}.
\]
\end{xsolution}

\subsection*{练习题 13}
\begin{xsolution}
先证必要性。若
\[
  z=f(x,y)=\varphi(ax+by),
\]
则
\[
  z_x=a\varphi'(ax+by),
  \qquad
  z_y=b\varphi'(ax+by),
\]
所以
\[
  bz_x=az_y.
\]

再证充分性。设
\[
  bz_x=az_y.
\]
作可逆线性变换
\[
  s=ax+by,
  \qquad
  t=bx-ay.
\]
其逆变换为
\[
  x=\frac{as+bt}{a^2+b^2},
  \qquad
  y=\frac{bs-at}{a^2+b^2}.
\]
令
\[
  Z(s,t)=z(x(s,t),y(s,t)).
\]
则
\begin{align*}
  Z_t
  &=z_x\frac{b}{a^2+b^2}
    +z_y\frac{-a}{a^2+b^2}\\
  &=\frac{bz_x-az_y}{a^2+b^2}=0.
\end{align*}
因此 $Z$ 与 $t$ 无关，存在一元函数 $\varphi$ 使
\[
  Z(s,t)=\varphi(s).
\]
还原变量得到
\[
  \boxed{z=f(x,y)=\varphi(ax+by)}.
\]
故所给条件充分且必要。
\end{xsolution}

\subsection*{练习题 14}
\begin{xsolution}
由
\[
  F(u_x,u_y)=0
\]
分别对 $x,y$ 求偏导，得到
\[
  F_su_{xx}+F_tu_{xy}=0,
\]
\[
  F_su_{xy}+F_tu_{yy}=0.
\]
即
\[
  \begin{pmatrix}
    u_{xx}&u_{xy}\\
    u_{xy}&u_{yy}
  \end{pmatrix}
  \binom{F_s}{F_t}=0.
\]
题设
\[
  F_s^2+F_t^2\ne0
\]
说明向量 $(F_s,F_t)^T$ 非零。于是上面的 $2\times2$ 矩阵有非零核向量，故其行列式为零：
\[
  \boxed{u_{xx}u_{yy}-u_{xy}^2=0}.
\]
\end{xsolution}

\subsection*{练习题 15}
\begin{xsolution}
令
\[
  r=\sqrt{x^2+y^2},
  \qquad u(x,y)=f(r).
\]
当 $r>0$ 时
\[
  u_x=f'(r)\frac xr,
  \qquad
  u_y=f'(r)\frac yr.
\]
计算得
\[
  u_{xx}+u_{yy}=f''(r)+\frac1r f'(r).
\]
故调和方程等价于
\[
  f''(r)+\frac1r f'(r)=0,
\]
即
\[
  (r f'(r))'=0.
\]
因此
\[
  r f'(r)=C_1,
  \qquad
  f(r)=C_1\ln r+C_2.
\]
所以在不含原点的区域上，径向解为
\[
  \boxed{u(x,y)=C_1\ln\sqrt{x^2+y^2}+C_2}.
\]
若题设的“有连续二阶偏导数”是指 $u$ 在包含原点的整个邻域内为 $C^2$ 函数，则 $\ln r$ 在 $r=0$ 无法连续延拓，必须有 $C_1=0$。此时
\[
  \boxed{u(x,y)=C_2},
\]
即只有常值解。
\end{xsolution}

\section*{19.4.2　练习题}

\subsection*{练习题 1}
\begin{xsolution}
题设给出
\[
  f(x)-f(x_0)=A(x-x_0)+\littleo(x-x_0),
\]
其中
\[
  \frac{|\littleo(x-x_0)|}{|x-x_0|}\to0
  \qquad(x\to x_0).
\]
由于 $h\mapsto Ah$ 是线性映射，这正是 $f$ 在 $x_0$ 处可微的定义，所以 $f$ 在 $x_0$ 可微，其全导数就是线性映射 $A$。

再直接核对 Jacobi 矩阵。令 $e_j$ 为第 $j$ 个标准基向量，取
\[
  x=x_0+t e_j.
\]
则
\[
  \frac{f(x_0+t e_j)-f(x_0)}t
  =Ae_j+\frac{\littleo(t e_j)}t.
\]
令 $t\to0$，得到第 $j$ 个偏导向量
\[
  \frac{\partial f}{\partial x_j}(x_0)=Ae_j.
\]
因此 Jacobi 矩阵的第 $j$ 列正是 $A$ 的第 $j$ 列，故
\[
  \boxed{Jf(x_0)=A}.
\]
\end{xsolution}

\subsection*{练习题 2}
\begin{xsolution}
由 $|f(t)|=1$，有
\[
  f(t)\cdot f(t)=1.
\]
两边对 $t$ 求导，得到
\[
  2f'(t)\cdot f(t)=0,
\]
所以
\[
  \boxed{f'(t)\cdot f(t)=0}.
\]

几何上，曲线 $f(t)$ 始终位于单位球面 $S^2$ 上，$f(t)$ 是从原点指向球面点的半径向量，而 $f'(t)$ 是曲线在该点的切向量。上式说明球面曲线的切向量与半径向量正交，即 $f'(t)$ 位于单位球面在 $f(t)$ 处的切平面内。
\end{xsolution}

\subsection*{练习题 3}
\begin{xsolution}
记向量值函数
\[
  F(x,y)=(u(x,y),v(x,y)).
\]
题设不等式可写为
\[
  |F(p)-F(q)|^2\ge C|p-q|^2,
  \qquad p,q\in\RR^2.
\]
固定 $p=(x,y)$，任取向量 $h\in\RR^2$，并令 $q=p+th$。则
\[
  \left|\frac{F(p+th)-F(p)}t\right|
  \ge\sqrt C\,|h|.
\]
令 $t\to0$。由于 $F$ 可微，左端趋于 $|JF(p)h|$，因此
\[
  |JF(p)h|\ge\sqrt C\,|h|,
  \qquad h\in\RR^2.
\]
若 $JF(p)h=0$，则上式强迫 $h=0$。所以 $JF(p)$ 的核只有零向量。它是 $2\times2$ 方阵，故可逆，于是
\[
  \boxed{\det JF(p)
  =\frac{\partial(u,v)}{\partial(x,y)}(p)\ne0}.
\]
由于 $p$ 任意，结论在 $\RR^2$ 上处处成立。
\end{xsolution}

\section*{19.5.2　参考题}

\subsection*{参考题 1}
\begin{xsolution}
记
\[
  s=\frac{x^2}{t}\ge0.
\]
由
\[
  u(x,t)=\frac1{2\sqrt\pi}t^{-1/2}\ee^{-x^2/(4t)}
\]
计算得
\[
  u_x=-\frac{x}{4\sqrt\pi}t^{-3/2}\ee^{-x^2/(4t)},
\]
\[
  u_{xx}
  =\frac1{2\sqrt\pi}t^{-3/2}\ee^{-s/4}
    \left(\frac{s}{4}-\frac12\right).
\]
因此
\[
  |u_{xx}|
  =\frac1{2\sqrt\pi}t^{-3/2}
    \left|\frac{s}{4}-\frac12\right|\ee^{-s/4}.
\]
给定 $0<\lambda<\frac14$，令
\[
  \mu=\frac14-\lambda>0.
\]
则
\[
  \left|\frac{s}{4}-\frac12\right|\ee^{-s/4}
  =\left|\frac{s}{4}-\frac12\right|\ee^{-\mu s}\ee^{-\lambda s}.
\]
函数
\[
  \left|\frac{s}{4}-\frac12\right|\ee^{-\mu s}
\]
在 $[0,+\infty)$ 上有界。设其上界为 $M_\lambda$，取
\[
  C=\frac{M_\lambda}{2\sqrt\pi},
\]
便有
\[
  \boxed{
  \left|\frac{\partial^2u}{\partial x^2}\right|
  \le Ct^{-3/2}\ee^{-\lambda x^2/t}}
\]
对一切 $t>0$ 成立。
\end{xsolution}

\subsection*{参考题 2}
\begin{xsolution}
记
\[
  X=x-\xi,
  \quad Y=y-\eta,
  \quad Z=z-\zeta,
  \quad R=\sqrt{X^2+Y^2+Z^2}.
\]
则在 $(x,y,z)\ne(\xi,\eta,\zeta)$，即 $R>0$ 时，
\[
  \Gamma=-\frac1{4\pi}R^{-1}.
\]
直接求导得
\[
  \Gamma_z=\frac1{4\pi}ZR^{-3}.
\]
因为 $|Z|\le R$，所以
\[
  |\Gamma_z|
  \le\frac1{4\pi}R^{-2}
  =\frac1{4\pi}(X^2+Y^2+Z^2)^{-1}.
\]

再求导：
\[
  \Gamma_{yz}
  =-\frac3{4\pi}YZR^{-5}.
\]
由 $|YZ|\le R^2$，有
\[
  |\Gamma_{yz}|
  \le\frac3{4\pi}R^{-3}
  =\frac3{4\pi}(X^2+Y^2+Z^2)^{-3/2}.
\]

第三次求导：
\[
  \Gamma_{xyz}
  =\frac{15}{4\pi}XYZR^{-7}.
\]
利用 $|XYZ|\le R^3$，得到
\[
  |\Gamma_{xyz}|
  \le\frac{15}{4\pi}R^{-4}
  =\frac{15}{4\pi}(X^2+Y^2+Z^2)^{-2}.
\]
故第三个不等式可取
\[
  \boxed{C=\frac{15}{4\pi}}.
\]
三式均得证。
\end{xsolution}

\subsection*{参考题 3}
\begin{xsolution}
由于 $u$ 无零点，在每个连通分支上 $u$ 的符号固定，可以定义
\[
  v(x,y)=\ln|u(x,y)|.
\]
计算
\[
  v_{xy}
  =\frac{u u_{xy}-u_xu_y}{u^2}.
\]
因此方程
\[
  uu_{xy}=u_xu_y
\]
等价于
\[
  v_{xy}=0.
\]

先证必要方向。若 $v_{xy}=0$，则 $v_x$ 与 $y$ 无关，所以存在一元函数 $a(x)$ 使
\[
  v_x(x,y)=a(x).
\]
对 $x$ 积分，得到
\[
  v(x,y)=A(x)+B(y).
\]
于是
\[
  |u(x,y)|=\ee^{A(x)}\ee^{B(y)}.
\]
把 $u$ 在该连通分支上的固定符号吸收到一个因子中，可写成
\[
  \boxed{u(x,y)=f(x)g(y)}.
\]

反之，若 $u=f(x)g(y)$，则
\[
  u_x=f'(x)g(y),
  \quad
  u_y=f(x)g'(y),
  \quad
  u_{xy}=f'(x)g'(y),
\]
从而
\[
  uu_{xy}
  =f(x)g(y)f'(x)g'(y)
  =u_xu_y.
\]
所以条件充分且必要。
\end{xsolution}

\subsection*{参考题 4：命题 19.3.1--19.3.4}
\begin{xsolution}
\textbf{命题 19.3.1。}

若 $f$ 是 $n$ 次齐次函数，则对 $t>0$
\[
  f(tx,ty)=t^n f(x,y).
\]
对 $t$ 求导并令 $t=1$，得到
\[
  xf_x(x,y)+yf_y(x,y)=nf(x,y).
\]

反之，设
\[
  xf_x+yf_y=nf.
\]
固定 $(x,y)$，定义
\[
  \Phi(t)=t^{-n}f(tx,ty),
  \qquad t>0.
\]
则
\begin{align*}
  \Phi'(t)
  &=t^{-n-1}
    \left[t\bigl(xf_x(tx,ty)+yf_y(tx,ty)\bigr)-nf(tx,ty)\right]\\
  &=0.
\end{align*}
故 $\Phi(t)$ 为常数，$\Phi(t)=\Phi(1)=f(x,y)$，即
\[
  f(tx,ty)=t^n f(x,y).
\]
所以 $f$ 为 $n$ 次齐次函数。

\textbf{命题 19.3.2。}

由命题 19.3.1，
\[
  xf_x+yf_y=nf.
\]
分别对 $x,y$ 求偏导：
\[
  xf_{xx}+yf_{xy}=(n-1)f_x,
\]
\[
  xf_{xy}+yf_{yy}=(n-1)f_y.
\]
第一式乘 $x$，第二式乘 $y$，相加得
\begin{align*}
  x^2f_{xx}+2xyf_{xy}+y^2f_{yy}
  &=(n-1)(xf_x+yf_y)\\
  &=n(n-1)f.
\end{align*}
这正是
\[
  \left(x\frac\partial{\partial x}
  +y\frac\partial{\partial y}\right)^2f
  =n(n-1)f
\]
按书中所说明的形式记号的含义展开后的等式。

\textbf{命题 19.3.3。}

由
\[
  f(tx,ty)=t^n f(x,y)
\]
对 $x$ 求偏导，注意左端链式法则给出因子 $t$：
\[
  t f_x(tx,ty)=t^n f_x(x,y).
\]
故
\[
  f_x(tx,ty)=t^{n-1}f_x(x,y).
\]
同理
\[
  f_y(tx,ty)=t^{n-1}f_y(x,y).
\]
所以 $f_x,f_y$ 都是 $(n-1)$ 次齐次函数。

\textbf{命题 19.3.4。}

在单位圆
\[
  S^1=\{(\xi,\eta):\xi^2+\eta^2=1\}
\]
上，$f$ 连续，所以
\[
  M=\max_{S^1}|f|<\infty.
\]
对任意 $(x,y)\ne(0,0)$，令
\[
  r=\sqrt{x^2+y^2},
  \qquad
  (\xi,\eta)=\left(\frac xr,\frac yr\right)\in S^1.
\]
由齐次性
\[
  f(x,y)=f(r\xi,r\eta)=r^n f(\xi,\eta).
\]
于是
\[
  |f(x,y)|\le M r^n
  =M(x^2+y^2)^{n/2}.
\]
取任意 $C\ge M$ 的正常数即得结论。
\end{xsolution}

\subsection*{参考题 5}
\begin{xsolution}
记所给 Vandermonde 型行列式为
\[
  u(x_1,\ldots,x_n).
\]

\textbf{(1)} 同时作平移 $x_j\mapsto x_j+t$。第 $k+1$ 行（$k=0,1,\ldots,n-1$）的元素变为
\[
  (x_j+t)^k
  =\sum_{\ell=0}^k\binom{k}{\ell}t^{k-\ell}x_j^\ell.
\]
因此新行列式的第 $k+1$ 行是原来第 $1$ 至第 $k+1$ 行的线性组合，并且原第 $k+1$ 行的系数为 $1$。这是一个对各行作的下三角线性变换，其对角元全为 $1$，故行列式不变：
\[
  u(x_1+t,\ldots,x_n+t)=u(x_1,\ldots,x_n).
\]
对 $t$ 求导并令 $t=0$，由链式法则得
\[
  \boxed{\sum_{i=1}^n\frac{\partial u}{\partial x_i}=0}.
\]

\textbf{(2)} 同时作伸缩 $x_j\mapsto t x_j$。第 $k+1$ 行提出因子 $t^k$，因此
\[
  u(tx_1,\ldots,tx_n)
  =t^{0+1+\cdots+(n-1)}u(x_1,\ldots,x_n)
  =t^{n(n-1)/2}u.
\]
所以 $u$ 是 $\frac{n(n-1)}2$ 次齐次函数。由 Euler 齐次函数公式
\[
  \boxed{
  \sum_{i=1}^n x_i\frac{\partial u}{\partial x_i}
  =\frac{n(n-1)}2u}.
\]
\end{xsolution}

\subsection*{参考题 6}
\begin{xsolution}
设
\[
  y=Ax,
  \qquad y_j=\sum_{i=1}^n a_{ji}x_i,
  \qquad F(x)=f(y).
\]

\textbf{(1)} 链式法则给出
\[
  \frac{\partial F}{\partial x_i}
  =\sum_{j=1}^n\frac{\partial f}{\partial y_j}a_{ji}.
\]
以列向量记梯度，则
\[
  \nabla_xF=A^T\nabla_yf.
\]
由于 $A$ 为正交矩阵，$AA^T=A^TA=I$，故
\begin{align*}
  \sum_{i=1}^n\left(\frac{\partial F}{\partial x_i}\right)^2
  &=|\nabla_xF|^2\\
  &=|A^T\nabla_yf|^2\\
  &=|\nabla_yf|^2\\
  &=\sum_{j=1}^n\left(\frac{\partial f}{\partial y_j}\right)^2.
\end{align*}

\textbf{(2)} 再求一次偏导：
\[
  \frac{\partial^2F}{\partial x_i^2}
  =\sum_{j,k=1}^n
  \frac{\partial^2f}{\partial y_j\partial y_k}a_{ji}a_{ki}.
\]
对 $i$ 求和，利用
\[
  \sum_{i=1}^n a_{ji}a_{ki}=\delta_{jk},
\]
得到
\begin{align*}
  \sum_{i=1}^n\frac{\partial^2F}{\partial x_i^2}
  &=\sum_{j,k=1}^n
    \frac{\partial^2f}{\partial y_j\partial y_k}
    \sum_{i=1}^n a_{ji}a_{ki}\\
  &=\sum_{j=1}^n\frac{\partial^2f}{\partial y_j^2}.
\end{align*}
两式均得证。
\end{xsolution}

\subsection*{参考题 7}
\begin{xsolution}
\textbf{(1)} 有
\[
  \frac{\partial(x_1,x_2)}{\partial(r,\theta)}
  =
  \begin{vmatrix}
    \cos\theta&-r\sin\theta\\
    \sin\theta&r\cos\theta
  \end{vmatrix}
  =\boxed{r}.
\]

\textbf{(2)} Jacobi 行列式为
\[
  \begin{vmatrix}
    \cos\theta_1&-r\sin\theta_1&0\\
    \sin\theta_1\cos\theta_2&r\cos\theta_1\cos\theta_2&-r\sin\theta_1\sin\theta_2\\
    \sin\theta_1\sin\theta_2&r\cos\theta_1\sin\theta_2&r\sin\theta_1\cos\theta_2
  \end{vmatrix}.
\]
从第二、三列分别提出 $r,r\sin\theta_1$，剩余行列式为 $1$，故
\[
  \boxed{
  \frac{\partial(x_1,x_2,x_3)}{\partial(r,\theta_1,\theta_2)}
  =r^2\sin\theta_1}.
\]

\textbf{(3)} 记 $J_m$ 为 $m$ 维变换的 Jacobi 行列式。令
\[
  \rho=r\sin\theta_1.
\]
则
\[
  x_1=r\cos\theta_1,
\]
而 $(x_2,\ldots,x_m)$ 恰是以半径 $\rho$ 和角变量 $\theta_2,\ldots,\theta_{m-1}$ 表示的 $(m-1)$ 维同型变换。又
\[
  \frac{\partial(x_1,\rho)}{\partial(r,\theta_1)}
  =
  \begin{vmatrix}
    \cos\theta_1&-r\sin\theta_1\\
    \sin\theta_1&r\cos\theta_1
  \end{vmatrix}
  =r.
\]
因此有递推关系
\[
  J_m(r,\theta_1,\ldots,\theta_{m-1})
  =r\,J_{m-1}(r\sin\theta_1,\theta_2,\ldots,\theta_{m-1}).
\]
由 $J_2=r$ 出发归纳，得到
\[
  \boxed{
  J_m
  =r^{m-1}
   \sin^{m-2}\theta_1
   \sin^{m-3}\theta_2\cdots
   \sin\theta_{m-2}}.
\]
这就是所求的 $m$ 维球坐标 Jacobi 行列式。
\end{xsolution}

\subsection*{参考题 8}
\begin{xsolution}
\textbf{(1)} 设
\[
  P=f_1f_2\cdots f_n,
\]
各因子均不为零。全微分的乘积法则给出
\[
  \dd P
  =\sum_{i=1}^n
    f_1\cdots f_{i-1}(\dd f_i)f_{i+1}\cdots f_n.
\]
两边除以 $P$，得
\[
  \frac{\dd P}{P}
  =\sum_{i=1}^n\frac{\dd f_i}{f_i}.
\]
若 $\delta_i$ 表示第 $i$ 个因子的有符号相对误差，则由
\[
  \Delta P\approx\dd P
\]
得到
\[
  \boxed{
  \delta(P)\approx\sum_{i=1}^n\delta_i}.
\]
若“相对误差”取绝对值，则相应有估计
\[
  |\delta(P)|\lesssim\sum_{i=1}^n|\delta_i|.
\]

\textbf{(2)} 由
\[
  \dd\ln|P|=\frac{\dd P}{P}
\]
以及
\[
  \ln|P|=\sum_{i=1}^n\ln|f_i|
\]
立即再次得到
\[
  \frac{\dd P}{P}=\sum_{i=1}^n\frac{\dd f_i}{f_i}.
\]

更一般地，设
\[
  Q=\frac{f_1f_2\cdots f_n}{g_1g_2\cdots g_n},
\]
且所有因子均不为零。则
\[
  \ln|Q|
  =\sum_{i=1}^n\ln|f_i|-
    \sum_{i=1}^n\ln|g_i|,
\]
所以
\[
  \frac{\dd Q}{Q}
  =\sum_{i=1}^n\frac{\dd f_i}{f_i}
   -\sum_{i=1}^n\frac{\dd g_i}{g_i}.
\]
因此有符号相对误差近似满足
\[
  \boxed{
  \delta(Q)
  \approx\sum_{i=1}^n\delta(f_i)
  -\sum_{i=1}^n\delta(g_i)}.
\]
若只估计误差大小，则
\[
  \boxed{
  |\delta(Q)|
  \lesssim\sum_{i=1}^n|\delta(f_i)|
  +\sum_{i=1}^n|\delta(g_i)|}.
\]
这就是通常所说的乘除运算相对误差相加的规则。
\end{xsolution}

\subsection*{参考题 9}
\begin{xsolution}
答案是：\textbf{必定线性相关}。

把各函数组成列向量
\[
  X(t)=(x_1(t),\ldots,x_n(t))^T.
\]
题设方程组写成
\[
  X'(t)=AX(t),
  \qquad A=(a_{ij}).
\]
于是对任意 $k\ge0$，
\[
  X^{(k)}(t)=A^kX(t).
\]
取某个固定 $t_0$，定义矩阵
\[
  K(t)=\bigl[X(t),X'(t),\ldots,X^{(n-1)}(t)\bigr].
\]
因为
\[
  X(t)=\ee^{A(t-t_0)}X(t_0)
\]
且 $A$ 与 $\ee^{A(t-t_0)}$ 可交换，所以
\[
  K(t)=\ee^{A(t-t_0)}K(t_0).
\]
从而
\[
  \det K(t)
  =\det\bigl(\ee^{A(t-t_0)}\bigr)\det K(t_0)
  =\ee^{\operatorname{tr}(A)(t-t_0)}\det K(t_0).
\]
题设 $a_{ij}\ge0$，特别地 $a_{ii}\ge0$，故
\[
  \operatorname{tr}(A)\ge0.
\]
另一方面，$X(t)\to0$，所以对每个固定 $k$，
\[
  X^{(k)}(t)=A^kX(t)\to0.
\]
于是 $K(t)$ 的每一列都趋于零，故
\[
  \det K(t)\to0.
\]
若 $\det K(t_0)\ne0$，上面的指数公式在 $\operatorname{tr}(A)\ge0$ 时不可能趋于零，矛盾。因此
\[
  \det K(t_0)=0.
\]

所以存在非零行向量 $c^T=(c_1,\ldots,c_n)$，使
\[
  c^TA^kX(t_0)=0,
  \qquad k=0,1,\ldots,n-1.
\]
由 Cayley--Hamilton 定理，$A^k$ 对所有 $k\ge n$ 都可表示成 $I,A,\ldots,A^{n-1}$ 的线性组合，故实际上
\[
  c^TA^kX(t_0)=0,
  \qquad k=0,1,2,\ldots.
\]
因此
\begin{align*}
  c^TX(t)
  &=c^T\ee^{A(t-t_0)}X(t_0)\\
  &=\sum_{k=0}^\infty
    \frac{(t-t_0)^k}{k!}c^TA^kX(t_0)=0.
\end{align*}
也就是
\[
  c_1x_1(t)+\cdots+c_nx_n(t)\equiv0,
\]
其中 $(c_1,\ldots,c_n)\ne0$。故 $x_1(t),\ldots,x_n(t)$ 线性相关。
\end{xsolution}

\subsection*{参考题 10}
\begin{xsolution}
记
\[
  B_{ij}(x)=\frac{\partial f_j}{\partial x_i}(x),
\]
则 $Jf=(B_{ij})$。对固定的 $i,j$，其代数余子式可以用 Levi--Civita 符号写为
\[
  A_{ij}
  =\frac1{(n-1)!}
  \sum_{i_2,\ldots,i_n}
  \sum_{j_2,\ldots,j_n}
  \varepsilon_{i i_2\cdots i_n}
  \varepsilon_{j j_2\cdots j_n}
  \prod_{k=2}^n
  \frac{\partial f_{j_k}}{\partial x_{i_k}}.
\]
这里各指标均从 $1$ 到 $n$ 求和；Levi--Civita 符号自动使有重复指标的项为零。

固定 $j$，对 $x_i$ 求导并对 $i$ 求和。乘积求导得到
\begin{align*}
  \sum_{i=1}^n\frac{\partial A_{ij}}{\partial x_i}
  &=\frac1{(n-1)!}
  \sum_{s=2}^n
  \sum_{i,i_2,\ldots,i_n}
  \sum_{j_2,\ldots,j_n}
  \varepsilon_{i i_2\cdots i_n}
  \varepsilon_{j j_2\cdots j_n}\\
  &\quad\times
  \frac{\partial^2 f_{j_s}}
       {\partial x_i\partial x_{i_s}}
  \prod_{\substack{k=2\\k\ne s}}^n
  \frac{\partial f_{j_k}}{\partial x_{i_k}}.
\end{align*}

现在固定 $s$ 以及其余所有指标，只考察关于 $i$ 与 $i_s$ 的求和。由于 $f\in C^2$，混合偏导可交换：
\[
  \frac{\partial^2 f_{j_s}}
       {\partial x_i\partial x_{i_s}}
  =
  \frac{\partial^2 f_{j_s}}
       {\partial x_{i_s}\partial x_i}.
\]
因此这一二阶偏导因子在交换 $i$ 与 $i_s$ 时不变；而
\[
  \varepsilon_{i i_2\cdots i_s\cdots i_n}
  =-
  \varepsilon_{i_s i_2\cdots i\cdots i_n}
\]
在同一交换下变号。其余乘积因子不受影响。故每一项都与交换 $i,i_s$ 后的项成对抵消；当 $i=i_s$ 时 Levi--Civita 符号本来就是零。

所以对每个 $s$ 的总和均为零，最终得到
\[
  \boxed{
  \sum_{i=1}^n
  \frac{\partial A_{ij}}{\partial x_i}(x)=0,
  \qquad j=1,2,\ldots,n}.
\]
这就是 Hadamard 恒等式。
\end{xsolution}
